% Source: https://artofproblemsolving.com/community/c3h2549840_rational_inequality_graph
% Stable identity: aops-topic-2549840-q01
$R$ is defined for $x\ne\pm2$. It has a zero at $x=1$, vertical asymptotes $x=-2$ and $x=2$, and $R(x)\to0$ as $x\to\pm\infty$ (horizontal asymptote $y=0$); also $R(0)=\dfrac14$.

The sign of $R(x)=\dfrac{x-1}{(x-2)(x+2)}$ changes only at $-2$, $1$ and $2$:
\[\begin{array}{c|cccc}
&x<-2&-2<x<1&1<x<2&x>2\\\hline
x-1&-&-&+&+\\
x^2-4&+&-&-&+\\
R(x)&-&+&-&+
\end{array}\]
Including the zero at $x=1$ (but not $\pm2$, where $R$ is undefined),
\[R(x)\ge0\iff\boxed{-2<x\le1\ \text{ or }\ x>2}.\]
In particular $R(x)\ge0$ does \emph{not} hold for every $x$: for example $R(-3)=-\dfrac45$.

In the sketch, the solution set is marked in red on the $x$-axis.

\begin{asy}
usepackage("amsmath");
size(7cm);
defaultpen(fontsize(10pt));
import graph;
real f(real x){return (x-1)/(x^2-4);}
draw((-5.3,0)--(5.4,0),Arrow(TeXHead));draw((0,-3.3)--(0,3.4),Arrow(TeXHead));
draw((-2,-3.2)--(-2,3.2),dashed);draw((2,-3.2)--(2,3.2),dashed);
draw(graph(f,-5,-2.255),blue+linewidth(0.8));
draw(graph(f,-1.755,1.921),blue+linewidth(0.8));
draw(graph(f,2.089,5),blue+linewidth(0.8));
draw((-2,0)--(1,0),red+linewidth(2));draw((2,0)--(5.3,0),red+linewidth(2));
filldraw(circle((-2,0),0.09),white,red);filldraw(circle((2,0),0.09),white,red);dot((1,0),red);
label("$x$",(5.4,0),E);label("$y$",(0,3.4),N);
label("$-2$",(-2,0),SW,fontsize(8pt));label("$2$",(2,0),SE,fontsize(8pt));label("$1$",(1,0),S,fontsize(8pt));
label("$x=-2$",(-2,3.2),N,fontsize(8pt));label("$x=2$",(2,3.2),N,fontsize(8pt));
\end{asy}
